% Generated draft: requires mathematical and source review.
\documentclass[11pt]{article}
\usepackage[margin=25mm]{geometry}
\usepackage{fontspec}
\setmainfont{Latin Modern Roman}
\usepackage{amsmath,amssymb,mathtools,unicode-math}
\setmathfont{Latin Modern Math}
\usepackage{enumitem,needspace,xurl,hyperref}
\setlength{\parindent}{0pt}
\setlength{\parskip}{5pt}
\setlength{\emergencystretch}{3em}
\newcommand{\N}{\mathbb{N}}
\newcommand{\Z}{\mathbb{Z}}
\newcommand{\Q}{\mathbb{Q}}
\newcommand{\R}{\mathbb{R}}
\newcommand{\C}{\mathbb{C}}
\newcommand{\F}{\mathbb{F}}
\newcommand{\A}{\mathbb{A}}
\newcommand{\PP}{\mathbb{P}}
\newcommand{\E}{\mathbb{E}}
\newcommand{\eps}{\varepsilon}
\newcommand{\dd}{\,\mathrm{d}}
\begin{document}
\section*{UnsolvedMath ID: 30001646}
\textbf{Source problem code:} catalogue entry
\subsection{Definitions and mathematical statement}
The supplied record concerns a proposed Brumer--Stark variant for a \(T\)-class group. It does not supply enough hypotheses or definitions for a self-contained conjecture. The following retains its formula and records its intended annihilation assertion conditionally.

Let \(S\) and \(T\) be disjoint finite sets of places in the unspecified arithmetic setting of the source. The displayed expression in the record is
\[
\theta_{S,T}(s)
=
\left(\prod_{P\in T}
\bigl(1-\sigma_P^{-1}(NP)^{1-s}\bigr)^{-1}\right)
\theta_S(s).
\]
The record does not specify the underlying field or extension, the group containing \(\sigma_P\), the definition and normalization of \(\theta_S(s)\), or the meaning of \(NP\). Nor does it specify the algebra in which the displayed inverses and the value at \(s=0\) are defined. The inverse Euler factors above are reproduced exactly as supplied.

Write \(A=\operatorname{Cl}(\mathcal O_T)\) for the object called the \(T\)-class group in the record. Its precise definition is not supplied; in particular, the title's reference to ray class groups does not determine its modulus or local conditions. For an element \(\alpha\) of a ring acting on an abelian group \(A\), to say that \(\alpha\) annihilates \(A\) means
\[
\alpha\cdot a=0\qquad\text{for every }a\in A.
\]
Thus the proposed assertion, whenever its arithmetic data, acting ring, action, and special value are defined, is
\[
\theta_{S,T}(0)\cdot a=0
\qquad\text{for every }a\in\operatorname{Cl}(\mathcal O_T).
\]
This assertion is subject to an unspecified ``minor technical condition on \(T\).'' The supplied source-verification note reports that the cited report does not state that condition. Consequently, the assertion cannot be quantified over all disjoint finite \(S,T\), and no precise admissible family can be extracted from the supplied record. Additional restrictions on \(S\), including any requirements concerning archimedean or ramified places, are also not supplied. This is a conditional transcription of the intended assertion, not a certified self-contained conjecture.
\subsection{Short English statement}
The proposed variant says that a modified special value \(\theta_{S,T}(0)\) kills every element of a \(T\)-class group. The underlying arithmetic setting, several definitions, and a required condition on \(T\) are missing, so the supplied record does not determine a precise unconditional formulation. Its supplied status is partially solved; this draft does not independently verify that assessment.
\subsection{Sources}
\begin{itemize}
\item Supplied local UnsolvedMath registry record, permanent ID \(30001646\), exact source problem code \texttt{catalogue entry}, title ``Brumer--Stark Conjecture for Ray Class Groups.'' The proposed year is recorded as 2010; no individual proposer is supplied.
\item Supplied source citation and literature-source label ``original Oberwolfach problem report'': \emph{Mini-Workshop: 1-Motives (2011)}, Oberwolfach Reports, pp.~3237--3267, DOI \texttt{10.4171/owr/2010/55}. \url{https://doi.org/10.4171/owr/2010/55}. The background attributes the collection to workshop open-problem discussions documented in Oberwolfach Reports. The citation is reproduced as supplied, including its year designation.
\item Supplied literature source, ``Brumer-Stark away from 2'': \url{https://annals.math.princeton.edu/2023/197-1/p05}.
\item Supplied literature source, ``equivariant Iwasawa applications'': \url{https://doi.org/10.1007/s40687-025-00578-8}.
\item Literature assessment dated exactly \texttt{2026-08-21}: the supplied curation labels the problem partially solved and states that the function-field analogue is a theorem, while the number-field \(T\)-class-group conjecture is known in substantial partial cases, including the odd-prime part. This is supplied dated curation derived from the OWR collection, not an independently verified literature assessment. Its time-sensitive status and cited claims require independent verification.
\item The supplied statement-verification field labels the formulation \texttt{unrecoverable}, reports that the official report leaves the technical condition on \(T\) unspecified, and says that no separate clean statement was promoted. This draft preserves that source limitation and does not claim to have inspected the report.
\end{itemize}
\end{document}
